Για τη διατύπωση του μαθηματικού προβλήματος ξεκινάμε με την καταγραφή των άγνωστων μεγεθών. Για την απλοποίηση της διατύπωσης, ξεκινάμε με τον υπολογισμό των γνωστών μεγεθών.\\
\noindent
Με βάση τις εξισώσεις που αναπτύχθηκαν στα \cref{sec:tyre_model,sec:vehicle_model}, το σύνολο των άγνωστων μεταβλητών δίνεται στον παρακάτω πίνακα.

\begin{table}[htbp]
    \centering
    \small
    \renewcommand{\arraystretch}{1.2}
    \begin{tabular}{@{}l l@{}}
        \toprule
        \textbf{Μεταβλητή} & \textbf{Περιγραφή} \\
        \midrule
        $\boldsymbol{\beta}$         & \textbf{Γωνία πλευρικής ολίσθησης οχήματος} \\
        $\boldsymbol{\delta}$        & \textbf{Γωνία διεύθυνσης εμπρός τροχών} \\
        \midrule
        $F_x$           & Συνολικό διαμήκες φορτίο οχήματος \\
        $F_y$           & Συνολικό πλευρικό φορτίο οχήματος \\
        \midrule
        $F_{x,y,z,FL}$  & Φορτία τροχού FL \\
        $F_{x,y,z,FR}$  & Φορτία τροχού FR \\
        $F_{x,y,z,RL}$  & Φορτία τροχού RL \\
        $F_{x,y,z,RR}$  & Φορτία τροχού RR \\
        \midrule
        $\kappa_{FL}$   & Λόγος ολίσθησης τροχού FL \\
        $\kappa_{FR}$   & Λόγος ολίσθησης τροχού FR \\
        $\kappa_{RL}$   & Λόγος ολίσθησης τροχού RL \\
        $\kappa_{RR}$   & Λόγος ολίσθησης τροχού RR \\
        \midrule
        $\alpha_{FL}$   & Γωνία ολίσθησης τροχού FL \\
        $\alpha_{FR}$   & Γωνία ολίσθησης τροχού FR \\
        $\alpha_{RL}$   & Γωνία ολίσθησης τροχού RL \\
        $\alpha_{RR}$   & Γωνία ολίσθησης τροχού RR \\
        \bottomrule
    \end{tabular}
    \caption{Σύνολο αγνώστων μοντέλου εκτίμησης κατάστασης.}
    \label{tab:stateUnknowns}
\end{table}

Από την κατάστρωση των αγνώστων προκύπτει πως το μοντέλο έχει 24 άγνωστα μεγέθη. Το σύστημα των 24 εξισώσεων δε διαθέτει λύση κλειστού τύπου επομένως θα επιλυθεί αριθμητικά.
Ωστόσο, η αριθμητική διαδικασία μπορεί να επιταχυνθεί αν πραγματοποιήσουμε μια προ-επεξεργασία, συσχετίζοντας αναλυτικά τις άγνωστες ποσότητες μεταξύ τους. Πράγματι, όπως θα αναπτυχθεί παρακάτω, όλα τα άγνωστα μεγέθη μπορούν να εκφρασθούν μονοσήμαντα συναρτήσει των γωνιών $\beta$ και $\delta$, μειώνοντας το πρόβλημα σε ένα σύστημα μη-γραμμικών εξισώσεων $2\times2$. Η διαδικασία των αντικαταστάσεων στις εξισώσεις καλύπτεται εν μέρει από το σχεδιάγραμμα της επίλυσης \ref{fig:slip_est_flow}. \\[6pt]

\subsubsection{Υπολογισμός ταχυτήτων και επιταχύνσεων}

\noindent
Αρχικά, υπολογίζονται οι επιταχύνσεις στο τοπικό σύστημα του οχήματος $\alpha_x, \alpha_y$ με μετασχηματισμό \ref{eq:vectorTransformations}.

\begin{equation}
    \begin{aligned}
         \begin{Bmatrix}
            \alpha_x \\ \alpha_y
        \end{Bmatrix}
        &= R(\beta)
        \begin{Bmatrix}
            \alpha_{\hat{X}} \\ \alpha_{\hat{Y}}
        \end{Bmatrix}\\[10pt]
         \begin{Bmatrix}
            U_x \\ U_y
        \end{Bmatrix}
        &= R(\beta)
        \begin{Bmatrix}
            U_{\hat{X}} \\ U_{\hat{Y}}
        \end{Bmatrix}
    \end{aligned}
    \label{eq:localAccelAndVelocity}
\end{equation}

\noindent
Έπειτα, υπολογίζονται οι συνιστώσες της ταχύτητας κάθε τροχού, μέσω της εξίσωσης (\ref{eq:wheelVelocity}).\\[5pt]

\noindent
Για τους εμπρόσθιους τροχούς, τα διανύσματα των ταχυτήτων μετασχηματίζονται στο τοπικό σύστημα των τροχών ως:

\begin{equation}
    \vec{U}_{uv,i} = R(-\delta)\vec{U}_{xy,i}
    \label{eq:wheelVelocityTransformation}
\end{equation}

 \noindent
 Επομένως, από την \cref{eq:slipAngleCalculation} υπολογίζονται οι γωνίες ολίσθησης κάθε τροχού.

\begin{equation}
    \alpha_i = tan^{-1}\Bigl(\dfrac{U_{v,i}}{U_{u,i}}\Bigr)
    \label{eq:slipAngleCalculation}
\end{equation}

\subsubsection{Υπολογισμός δυνάμεων}
\noindent
Έχοντας υπολογίσει τις επιταχύνσεις και τις ταχύτητες του οχήματος στο τοπικό του σύστημα αξόνων, υπολογίζονται τα συνολικά φορτία του αμαξώματος $F_x, F_y$, μέσω των εξισώσεων (\ref{eq:vehicleMotion1}) και (\ref{eq:vehicleMotion2}).

\noindent
Παράλληλα, υπολογίζονται μέσω των εξισώσεων (\ref{eq:drag})--(\ref{eq:liftRear}) η οπισθέλκουσα $F_{x,drag}$ και η αρνητική άνωση (downforce) στους δύο άξονες $F_{z,f,lift}, F_{z,r,lift}$.\\[5pt]

Και τα κατακόρυφα φορτία των τροχών $F_{z,i}$ υπολογίζονται από τις σχέσεις (\ref{eq:longitudinalWeightTransfer}), (\ref{eq:lateralWeightTransfer}) και (\ref{eq:normalLoads}). Τα τελικά κατακόρυφα φορτία διορθώνονται σε περίπτωση απώλειας επαφής, για αποφυγή λανθασμένων αποτελεσμάτων.

\begin{equation}
    \begin{aligned}
        F_{z,i,\text{διορθ.}} &= max(F_{z,i}, 0)\\
        F_{z,i,\text{παραβ.}} &= min(F_{z,i}, 0)
    \end{aligned}
    \label{eq:fzViolation}
\end{equation}

\noindent
Όπου $F_{z,i,\text{παραβ.}}$, είναι το ποσό του φορτίου που παραβιάζει τη συνθήκη επαφής.\\[2.5pt]

\noindent
Ένα αποτέλεσμα αυτής της προσέγγισης είναι πως όταν ένας τροχός χάσει επαφή με το έδαφος, δεν επανα-κατανέμεται το κατακόρυφο φορτίο στους υπόλοιπους τρεις τροχούς, επομένως δεν ισχύει πλέον ισορροπία δυνάμεων στον κατακόρυφο άξονα, ωστόσο, αυτό αναμένεται να μην εισάγει παρατηρήσιμο σφάλμα στα αποτελέσματα.\\[5pt]

Έχοντας υπολογίσει τις γωνίες ολίσθησης και τα κατακόρυφα φορτία, πλέον τα εγκάρσια φορτία υπολογίζονται από το μοντέλο ελαστικών της μορφής της (\ref{eq:tyreModelFinal}). Το μοντέλο ελαστικών λαμβάνει ως είσοδο το απαιτούμενο διαμήκες φορτίο, το οποίο κατά περίπτωση θα αναλυθεί στην επόμενη παράγραφο.\\[5pt]

\noindent
Τα φορτία των ελαστικών μετασχηματίζονται τελικά στο σύστημα του οχήματος:

\begin{equation}
    \vec{F}_{xy,i} = R(\delta)\vec{F}_{uv,i}
    \label{eq:tyreForcesTransformation}
\end{equation}

\subsubsection{Υπολογισμός διαμήκων φορτίων}
\noindent
Η συνισταμένη διαμήκης δύναμη του οχήματος δίνεται παρακάτω.

\begin{equation}
    F_x = F_{x,tires} + F_{x,drag}
    \label{eq:vehicleFx}    
\end{equation}

\noindent
Όπου $F_{x,tires}$ η συνισταμένη διαμήκης δύναμη που παράγουν τα ελαστικά, εκφρασμένη στο σύστημα αξόνων του οχήματος. Αντικαθιστώντας την (\ref{eq:vehicleFx}) στην \cref{eq:vehicleMotion1}, προκύπτει:

\begin{equation}
    F_{x,tires} = M\dfrac{d}{dt}U_x(t) - M\dot{\psi} U_y - F_{x,drag}
    \label{eq:modifiedVehicleMotionX}
\end{equation}

\noindent
Όπου $F_{x,tires}$
\begin{equation}
   F_{x,tires} = F_{x,FL} +  F_{x,FR} +F_{x,RL} +F_{x,RR}
   \label{eq:tyreTotalFx}
\end{equation}

\noindent
Όπου οι δύο πρώτοι όροι του δεξιού μέλους είναι ήδη γνωστοί από τις εξισώσεις (\ref{eq:localAccelAndVelocity}). Καθώς η τιμή του $F_{x,drag}$ είναι εξ'ορισμού αρνητική, η \cref{eq:modifiedVehicleMotionX} μας δείχνει πως όταν το όχημα επιβραδύνει, η οπισθέλκουσα μειώνει το απαιτούμενο φορτίο πέδησης που πρέπει να παρέχουν τα ελαστικά. Αντίθετα, όταν το όχημα επιταχύνει, τα ελαστικά πρέπει, εκτός από τα αδρανειακά, να παρέχουν επιπλέον φορτίο για να υπερνικήσουν την οπισθέλκουσα.\\[6pt]

\noindent
Επομένως, πλέον είναι καθορισμένο το συνολικό διαμήκες φορτίο που πρέπει να παράξουν οι τροχοί. Για τον καθορισμό της κατανομής του φορτίου στους τροχούς ($F_{x,i}, i\in \{FL, FR, RL, RR\}$), διακρίνουμε δύο περιπτώσεις, την πρόωση και την πέδηση. 
Αναπτύσσοντας το συνολικό φορτίο που παράγουν οι τροχοί, οι κινητήριοι τροχοί (πίσω τροχοί) λαμβάνουν το παρακάτω φορτίο.

\begin{equation}
    F_{x,rear} = F_{x,RL} + F_{x,RR} = F_{x,tires} - F_{x,FL} - F_{x,FR}
    \label{eq:rearTractiveForce}
\end{equation}

\begin{figure}[htbp]
    \centering
    % figures/fig_steering_and_forces.tex
% Wheel steered by angle delta with forces F_x, F_y in the local frame
\begin{tikzpicture}[>=Latex, thick, font=\small]

    \def\deltaFig{25}          % steering angle in degrees

    % ---- Wheel rectangle (top view), rotated by delta about the origin ----
    \begin{scope}[rotate=\deltaFig]
    \draw[thick, fill=black!5, rounded corners=8pt] (-1.5, -2.5) rectangle (1.5, 2.5);
    \end{scope}

    % ---- Vehicle reference frame axes ----
    \draw[dashed,->, thick] (0,0) -- (0,5) node[right] {$x$};
    \draw[dashed,->, thick] (0,0) -- (-4,0) node[below]  {$y$};

 % ---- Wheel direction (rotated by delta from vehicle x) ----
    \draw[dashed, gray!100, ->] (0,0) -- ({5*sin(-\deltaFig)}, {5*cos(-\deltaFig)}) node[right] {$u$};

    \draw[dashed, gray!100, ->] (0,0) -- ({-5*cos(-\deltaFig)}, {5*sin(-\deltaFig)}) node[below] {$v$};

    % ---- Reference dashed line (vehicle x direction, along wheel center) ----
    \draw[dashed, gray!60] (0,-2) -- (0,0);

    % ---- Steering angle arc ----
    \draw[<->] (0,4) arc[start angle=90, end angle={90+\deltaFig}, radius=4]
        node[midway, above, yshift=1pt] {$\delta$};

    % ---- Wheel-center coordinate (used to anchor force vectors) ----
    \coordinate (W) at ({0}, {0});
    \fill (W) circle (1pt);

    % ---- Forces in the WHEEL local frame ----
    % F_x along wheel longitudinal axis (rolling direction)
    % F_y perpendicular to wheel plane (lateral)
    \draw[->, thick] (W) -- ++({-3*cos(-\deltaFig)}, {3*sin(-\deltaFig)})
        node[above left]  {$F_v$};

    \draw[->, thick] (0,0) -- (0,{3*sin(-\deltaFig)}) node[right] {$F_x$};
    \draw[dashed, gray!60] (0,{3*sin(-\deltaFig)}) -- ({-3*cos(-\deltaFig)},{3*sin(-\deltaFig)});
    \draw[->, thick] (0,0) -- ({-3*cos(-\deltaFig)},0) node[above] {$F_y$};
    \draw[dashed, gray!60] ({-3*cos(-\deltaFig)},0) -- ({-3*cos(-\deltaFig)},{3*sin(-\deltaFig)});

\end{tikzpicture}
    \caption{Ανάλυση δυνάμεων κατευθυντήριου τροχού}
    \label{fig:inducedDrag}
\end{figure}

Στο \cref{fig:inducedDrag} παρατηρείται πως ένας κατευθυντήριος τροχός, ακόμη και αν δεν παράγει διαμήκη δύναμη στο τοπικό του σύστημα, εισάγει διαμήκες φορτίο στο σύστημα αξόνων του οχήματος. Αυτός είναι ένας από τους λόγους που όταν ένα όχημα στρίβει, δε διατηρεί την ταχύτητά του αλλά επιβραδύνει.

Επομένως, οι περιπτώσεις πρόωσης και πέδησης διακρίνονται παρακάτω. Ολόκληρη η διαδικασία υπολογισμού των διαμήκων φορτίων συνοψίζεται στο \cref{fig:force_distribution_flow}.

\[
    \begin{cases}
        \text{Πρόωση}   & \text{αν } F_{x,rear} \geq 0 \\
        \text{Πέδηση}   & \text{αν } F_{x,rear} < 0 
    \end{cases}
\]


\paragraph{Περίπτωση πρόωσης}
Στην περίπτωση που το όχημα βρίσκεται σε κατάσταση πρόωσης, οι κινητήριοι τροχοί (πίσω) παρέχουν ολόκληρο το ποσό του φορτίου πρόωσης, ενώ για τους εμπρόσθιους τροχούς ισχύει $F_{u,FL},F_{u,FR} = 0$. Η διανομή του συνολικού φορτίου πρόωσης στους δύο τροχούς του άξονα πραγματοποιείται μέσω του διαφορικού, και μοντελοποιείται μέσω της σχέσης (\ref{eq:differential}), αντικαθιστώντας την \cref{eq:rearTractiveForce}.

\paragraph{Περίπτωση πέδησης}
Στην κατάσταση πέδησης, τα φορτία διανέμονται στον εμπρός και πίσω άξονα μέσω της σχέσης (\ref{eq:brakingDistribution}). Τα φορτία του κάθε άξονα θεωρούνται πως ισο-μοιράζονται στους δύο τροχούς.

\paragraph{Κορεσμός}
Σε δυναμικές συνθήκες, το ελαστικό μπορεί να μην έχει τη δυνατότητα να παράξει το απαιτούμενο φορτίο στις συνθήκες που βρίσκεται. Σε αυτή την περίπτωση, το διάνυσμα της λύσης του συστήματος κρίνεται ανέφικτο, αποδίδεται στο μοντέλο το μέγιστο φορτίο που είναι διαθέσιμο από το ελαστικό και καταγράφεται το έλλειμμα για τους σκοπούς των περιορισμών που αναλύεται σε επόμενη παράγραφο. 

\begin{equation}
    F_{x,i,\text{έλλειμμα}} = max\Bigl(
    \bigl(F_{x,i,\text{ζητ.}} - F_{x,i,\text{διαθ.}}\bigr), 0\Bigr)
    \label{eq:fxViolation}
\end{equation}

\begin{figure}[htbp!]
\centering
\begin{tikzpicture}[
    node distance=0.6cm and 1.0cm,
    block/.style   ={rectangle, draw=black, fill=white,   text width=8cm,
                     align=center, rounded corners=4pt, minimum height=0.9cm, font=\small},
    io/.style      ={rectangle, draw=black, fill=black!5, text width=6cm,
                     align=center, rounded corners=2pt, minimum height=0.9cm, font=\small},
    outbox/.style  ={rectangle, draw=black, fill=black!12, text width=6cm,
                     align=center, rounded corners=2pt, minimum height=0.9cm, font=\small\bfseries},
    decision/.style={diamond,   draw=black, fill=white,   text width=2.6cm,
                     align=center, font=\scriptsize, inner sep=1pt},
    arr/.style     ={->, thick, black}
]

  \node[block]          (fyFront)  {Υπολογισμός πλευρικής δύναμης εμπρόσθιων τροχών $F_{y,front,0}$\\ υπό υπόθεση $F_x = 0$};
  \node[block, below=of fyFront]                  (fxRear)   {Εκτίμηση απαιτούμενης διαμήκους δύναμης οπίσθιων τροχών ($F_{x,rear,0}$)};
  \node[decision, below=of fxRear]                (sign)     {$F_{x,rear,0} \geq 0$;};

  % --- Traction branch (right) ---
  \node[block, right=1.6cm of sign, xshift=1.2cm] (traction) {Πρόωση:\\ κατανομή $F_{x,rear,0}$  στους δύο\\ κινητήριους τροχούς};

  % --- Braking branch (left) ---
  \node[block, below=1.4cm of sign]               (splitBrk) {Πέδηση:\\ διανομή συνολικής δύναμης\\ πέδησης εμπρός/πίσω $F_{x,front}, F_{x,rear}$};
  \node[block, below=of splitBrk]                 (recFy)    {Επανυπολογισμός $F_{y,front}$ λαμβάνοντας υπόψη νέα $F_{x,front}$};
  \node[block, below=of recFy]                    (recFxr)   {Επανυπολογισμός $F_{x,rear}$ και κατανομή στους κινητήριους τροχούς};

  % --- Output ---
  \node[outbox, below=1.4cm of recFxr]            (out)      {Έξοδος: δυνάμεις $F_x, F_y$ ανά τροχό};

  % --- Arrows ---
  \draw[arr] (fyFront) -- (fxRear);
  \draw[arr] (fxRear)  -- (sign);

  % Braking branch (straight down)
  \draw[arr] (sign)    -- node[right]{όχι (πέδηση)} (splitBrk);
  \draw[arr] (splitBrk) -- (recFy);
  \draw[arr] (recFy)   -- (recFxr);
  \draw[arr] (recFxr)  -- (out);

  % Traction branch (right, then down and back)
  \draw[arr] (sign.east) -- node[above]{ναι (πρόωση)} (traction.west);
  \draw[arr] (traction.south) |- (out.east);

\end{tikzpicture}
\caption{Διάγραμμα ροής -- διαδικασία υπολογισμού διαμήκων φορτίων πρόωσης και πέδησης.}
\label{fig:force_distribution_flow}
\end{figure}

\subsubsection{Τελικό σύστημα εξισώσεων}

Μετά από όλους του υπολογισμούς του κεφαλαίου, το σύστημα εξισώσεων προς επίλυση διαμορφώνεται από τις δύο εξισώσεις κίνησης (\ref{eq:vehicleMotion2}) και (\ref{eq:vehicleMotion3}) σε μόνιμη κατάσταση $\Bigl(\dfrac{d}{dt}\dot{\psi}(t)\Bigr)$.

\begin{equation}
    \left\{
    \begin{aligned}
        M\dfrac{d}{dt}U_y(t) + M\dot{\psi}U_x - \sum_{i=FL,\dots,RR}\Bigl(F_{y,i}\Bigr) &= 0 \\[5pt]
        \sum_{i=FL,\dots,RR} \Bigl(\vec{r}_i\times\vec{F}_i\Bigr)_z &= 0
    \end{aligned}
    \right.
    \label{eq:stateSolverEquationSystem}
\end{equation}
